\newcounter{jessecount}
\setcounter{jessecount}{1}

\begin{list}{[\arabic{jessecount}]\addtocounter{jessecount}{1}}{\leftmargin=3em \itemsep=4pt}

\item
 Jesse Thaler and Ken Van Tilburg,
\emph{Identifying Boosted Objects with N-subjettiness},
JHEP 1103:015 (2011)
[arXiv:1011.2268].

\begin{quote}
This paper introduced $N$-subjettiness, a collider observable that probes a jet for $N$ prong substructure.  Jets are collimated sprays of particles that arise from the fragmentation of high-energy quarks and gluons.  Jets can also form if a heavy resonance---like a $W$ boson, $Z$ boson, or top quark---is produced with a large Lorentz boost.  Distinguishing between these scenarios is crucial in the search for new physics beyond the Standard Model.  $N$-subjettiness is one of the most powerful observables ever designed for this purpose, and it has been widely applied at the Large Hadron Collider.
\end{quote}
\item
 Andrew J.\ Larkoski, Simone Marzani, Gregory Soyez, and Jesse Thaler,
\emph{Soft Drop},
JHEP 1405:146 (2014)
[arXiv:1402.2657].

\begin{quote}
This paper introduced Soft Drop, a jet reclustering algorithm that removes soft (i.e. low energy) contimination.  In the complex environment of the Large Hadron Collider, there are multiple sources of contimination that polute the substructure of jets, including underlying event, initial state radiation, and pileup.  Soft drop---true to its name---is one of the most effective algorithms for mitigating soft radiation.  Moreover, soft drop preserves the hard (i.e. high energy) structure of jets, such that the properties of soft-dropped jets can still be predicted from first principles.
\end{quote}
\item
 Patrick T.\ Komiske, Eric M.\ Metodiev, and Jesse Thaler,
\emph{The Metric Space of Collider Events},
Phys.\ Rev.\ Lett.\ 123:041801 (2019)
[arXiv:1902.02346].

\begin{quote}
This paper showed how to define the geometry of collider events, based on techniques from optimal transport.  In optimal transport, one computes the amount of ``work" needed to transform one distribution into another, which yields a precise notion of distance.  In the collider context, one can consider distributions of energy deposited in an idealized detector at infinity.  Once one has defined the distance between any two collider events, one can compute many different geometric objects with wide applicability in collider physics.
\end{quote}
\item
 Patrick T.\ Komiske, Ian Moult, Jesse Thaler, and Hua Xing Zhu,
\emph{Analyzing N-point Energy Correlators Inside Jets with CMS Open Data},
Phys.\ Rev.\ Lett.\ 130:051901 (2023)
[arXiv:2201.07800].

\begin{quote}
This paper performed the first ever analysis of $N$-point correlators inside jets, using public data from the Large Hadron Collider.  $N$-point correlators (distinct from $N$-subjettiness) are powerful ways to probe the substructure of jets.  While $2$-point correlators had been studied in the 1970s, the $N$-point generalization was not pursued at the time.  This work not only showed the feasibility of $N$-point analyses, but it also revealed a clear phase transition between the partonic and hadronic phases of quantum chromodynamics, which inspired multiple follow-up studies in both particle and nuclear physics.
\end{quote}
\item
 Benoît Assi, Stefan Höche, Kyle Lee, and Jesse Thaler,
\emph{QCD Theory Meets Information Theory},
Phys.\ Rev.\ Lett.\ 135:131901 (2025)
[arXiv:2501.17219].

\begin{quote}
This paper introduced a new strategy to combine multiple first-principles calculations into a coherent theoretical prediction using ideas from information theory and machine learning.  The idea is to take a comprehensive but low fidelity prior prediction over all of phase space, and upgrade it using constraints from high fidelity precision calculations that target particular phase space configurations.  From a theoretical perspective, the key innovation is the use of logarithmic moment constraints, which have never before been studied in the literature but capture important information about the resummation regime of quantum chromodynamics.
\end{quote}
\end{list}
